\documentclass[11pt, a4paper]{article}

% bfw-style.tex — gemeinsame Style-Definitionen für alle Handouts/Aufgabenhefte
% Voraussetzung: Kompilation mit xelatex (wegen fontspec)

\usepackage[ngerman]{babel}
\usepackage{geometry}
\geometry{a4paper, top=2.2cm, bottom=2.2cm, left=2.3cm, right=2.3cm}

\usepackage{fontspec}
% Fallback-Kette: Source Sans Pro -> Calibri -> Segoe UI -> Latin Modern Sans
% (Latin Modern ist immer mit MiKTeX vorhanden)
\IfFontExistsTF{Source Sans Pro}{%
  \setmainfont{Source Sans Pro}[Ligatures=TeX]%
}{\IfFontExistsTF{Calibri}{%
  \setmainfont{Calibri}[Ligatures=TeX]%
}{\IfFontExistsTF{Segoe UI}{%
  \setmainfont{Segoe UI}[Ligatures=TeX]%
}{%
  \setmainfont{Latin Modern Sans}[Ligatures=TeX]%
}}}
\IfFontExistsTF{Source Code Pro}{%
  \setmonofont{Source Code Pro}[Scale=0.92]%
}{\IfFontExistsTF{Consolas}{%
  \setmonofont{Consolas}[Scale=0.92]%
}{%
  \setmonofont{Latin Modern Mono}[Scale=0.92]%
}}

\usepackage{xcolor}
% Farbpalette: hell, freundlich, modern
\definecolor{bfwPrimary}{HTML}{0EA5A4}    % Petrol/Türkis - Primär
\definecolor{bfwPrimaryDark}{HTML}{0F766E} % dunkler Petrol für Headings
\definecolor{bfwAccent}{HTML}{F59E0B}     % Amber - Akzent
\definecolor{bfwSuccess}{HTML}{10B981}    % Grün - Erfolg / Lösung
\definecolor{bfwWarn}{HTML}{F97316}       % Orange - Achtung
\definecolor{bfwInk}{HTML}{0F172A}        % fast Schwarz
\definecolor{bfwInkSoft}{HTML}{475569}    % gedämpftes Grau
\definecolor{bfwBgSoft}{HTML}{F8FAFC}     % off-white
\definecolor{bfwBgTeal}{HTML}{ECFEFF}     % helles Türkis
\definecolor{bfwBgAmber}{HTML}{FEF3C7}    % helles Gelb
\definecolor{bfwBgRose}{HTML}{FFE4E6}     % helles Rosé für Eselsbrücken

\usepackage[colorlinks=true, linkcolor=bfwPrimaryDark, urlcolor=bfwPrimaryDark]{hyperref}
\usepackage{enumitem}
\setlist[itemize]{leftmargin=*, itemsep=2pt, topsep=2pt}
\setlist[enumerate]{leftmargin=*, itemsep=2pt, topsep=2pt}

\usepackage[most]{tcolorbox}
\usepackage{booktabs}
\usepackage{tikz}
\usetikzlibrary{decorations.pathmorphing, positioning, shapes.geometric, arrows.meta}

% Merkkästchen — wichtige Definition / Kernpunkt
\newtcolorbox{merksatz}[1][]{%
  enhanced, breakable,
  colback=bfwBgTeal,
  colframe=bfwPrimary,
  boxrule=0.6pt,
  arc=4pt,
  left=10pt, right=10pt, top=8pt, bottom=8pt,
  fonttitle=\bfseries\color{bfwPrimaryDark},
  title={\faLightbulb\ Merksatz},
  #1
}

% Eselsbrücke — Mnemoniks
\newtcolorbox{eselsbruecke}[1][]{%
  enhanced, breakable,
  colback=bfwBgRose,
  colframe=bfwAccent,
  boxrule=0.6pt,
  arc=4pt,
  left=10pt, right=10pt, top=8pt, bottom=8pt,
  fonttitle=\bfseries\color{bfwWarn},
  title={Eselsbrücke},
  #1
}

% Beispiel-Box
\newtcolorbox{beispiel}[1][]{%
  enhanced, breakable,
  colback=bfwBgSoft,
  colframe=bfwInkSoft,
  boxrule=0.4pt,
  arc=3pt,
  left=10pt, right=10pt, top=8pt, bottom=8pt,
  fonttitle=\bfseries\color{bfwInk},
  title={Beispiel},
  #1
}

% Lösungs-Box (für Aufgabenhefte)
\newtcolorbox{loesung}[1][]{%
  enhanced, breakable,
  colback=bfwBgAmber!40,
  colframe=bfwSuccess,
  boxrule=0.6pt,
  arc=3pt,
  left=10pt, right=10pt, top=8pt, bottom=8pt,
  fonttitle=\bfseries\color{bfwSuccess!70!black},
  title={Lösung},
  #1
}

% Glossar-Eintrag
\newcommand{\glossareintrag}[2]{%
  \par\noindent\textbf{\color{bfwPrimaryDark}#1} \enspace #2\par\smallskip
}

% Section-Stil — etwas farbiger als Standard
\usepackage{titlesec}
\titleformat{\section}
  {\Large\bfseries\color{bfwPrimaryDark}}
  {\thesection}{0.6em}{}
\titleformat{\subsection}
  {\large\bfseries\color{bfwPrimary}}
  {\thesubsection}{0.5em}{}
\titleformat{\subsubsection}
  {\normalsize\bfseries\color{bfwInk}}
  {\thesubsubsection}{0.4em}{}

% TikZ-Stil "skizzenhaft" — etwas weicher als Rough.js, aber stabil
\tikzset{
  roughbox/.style = {
    draw=bfwPrimaryDark, thick, fill=bfwBgTeal,
    rounded corners=4pt, inner sep=6pt
  },
  roughaccent/.style = {
    draw=bfwAccent, thick, fill=bfwBgAmber,
    rounded corners=4pt, inner sep=6pt
  },
  roughline/.style = {
    draw=bfwInkSoft, thick, -{Stealth[length=2.5mm]}
  }
}

% Bullet-Symbole (bewusst kein FontAwesome — vermeidet Font-Installations-Probleme)
\newcommand{\faLightbulb}{$\bullet$}
\newcommand{\faBookmark}{$\bullet$}

% Header/Footer
\usepackage{fancyhdr}
\pagestyle{fancy}
\fancyhf{}
\renewcommand{\headrulewidth}{0pt}
\fancyfoot[L]{\small\color{bfwInkSoft}\thepage}
\fancyfoot[R]{\small\color{bfwInkSoft} BFW M-064 Beschaffung im Büromanagement}


\title{\vspace*{-1.5cm}\Huge\bfseries\color{bfwPrimaryDark}Aufgabenheft Tag~14\\[0.4em]
  \large\color{bfwInkSoft}\textnormal{Digitales DMS, Datenschutz \& Umweltschutz --- Übungen im IHK-Klausurformat}}
\author{\textsf{Büroprozesse gestalten und Arbeitsvorgänge organisieren \textperiodcentered{} M-067}\\
\textsf{\small\copyright\ Can Siebert\,\textperiodcentered\,2026}}
\date{Selbstlernmaterial \textperiodcentered{} 16.07.2026}

\begin{document}
\maketitle
\thispagestyle{empty}

\begin{merksatz}[title={\bfseries So nutzen Sie dieses Aufgabenheft}]
Bearbeiten Sie alle Aufgaben zunächst \emph{ohne} in die Lösungen zu schauen. Schreiben
Sie Ihre Antworten in vollständigen Sätzen --- die IHK-Prüfung verlangt formulierte
Begründungen, nicht bloße Stichworte. Beachten Sie die \emph{Operatoren} (nennen,
erläutern, beurteilen, begründen) --- sie geben den geforderten Antwort-Tiefegrad vor.
Erst danach öffnen Sie den Lösungsteil ab Seite~\pageref{loesungen} und vergleichen.
\end{merksatz}

\tableofcontents
\vspace{1em}
\hrule

\section{Aufgaben}

\subsection*{Aufgabe~1 --- Dokumentenmanagementsystem (12 Punkte)}

\noindent\textbf{\color{bfwAccent}YouTube-Vorbereitung:}\enspace \href{https://www.youtube.com/watch?v=D0oC16dJTdw}{DMS einfach erklärt} (\url{https://www.youtube.com/watch?v=D0oC16dJTdw})

\medskip
Die \emph{Duisdorfer BüroKonzept KG} hat ein hohes Schriftgutaufkommen und überlegt, ein
Dokumentenmanagementsystem (DMS) einzuführen.

\begin{enumerate}[label=(\alph*)]
  \item \textbf{Erläutern} Sie, was ein DMS ist und wie es arbeitet. \hfill (3~P.)
  \item \textbf{Nennen} Sie vier Vorteile eines DMS. \hfill (4~P.)
  \item \textbf{Erklären} Sie die vier Kernanforderungen der revisionssicheren Archivierung. \hfill (3~P.)
  \item \textbf{Beurteilen} Sie, welchen Beitrag ein DMS zum papierlosen Büro leistet. \hfill (2~P.)
\end{enumerate}

\subsection*{Aufgabe~2 --- Datenschutz vs.\ Datensicherheit (12 Punkte)}

\noindent\textbf{\color{bfwAccent}YouTube-Vorbereitung:}\enspace \href{https://www.youtube.com/watch?v=-QsJ5c-Fxbs}{Datenschutz und Datensicherheit} (\url{https://www.youtube.com/watch?v=-QsJ5c-Fxbs})

\medskip
\begin{enumerate}[label=(\alph*)]
  \item \textbf{Grenzen} Sie Datenschutz und Datensicherheit voneinander \textbf{ab} (Schutzobjekt und Gefahr). \hfill (4~P.)
  \item \textbf{Definieren} Sie den Begriff \emph{personenbezogene Daten} und \textbf{nennen} Sie drei Beispiele. \hfill (4~P.)
  \item \textbf{Nennen} Sie die beiden zentralen Rechtsgrundlagen des Datenschutzes. \hfill (2~P.)
  \item \textbf{Beurteilen} Sie die Aussage: \emph{„Eine Firewall ist eine Maßnahme des Datenschutzes.”} \hfill (2~P.)
\end{enumerate}

\subsection*{Aufgabe~3 --- Datenschutzbeauftragter \& Betroffenenrechte (14 Punkte)}

\noindent\textbf{\color{bfwAccent}YouTube-Vorbereitung:}\enspace \href{https://www.youtube.com/watch?v=4gRbl55hnSk}{Datenschutzbeauftragter \& DSGVO} (\url{https://www.youtube.com/watch?v=4gRbl55hnSk})

\medskip
Die \emph{Duisdorfer BüroKonzept KG} beschäftigt 25 Personen ständig mit der
automatisierten Verarbeitung personenbezogener Daten.

\begin{enumerate}[label=(\alph*)]
  \item \textbf{Begründen} Sie mit der gesetzlichen Schwelle, ob ein Datenschutzbeauftragter zu ernennen ist. \hfill (3~P.)
  \item \textbf{Nennen} Sie drei Aufgaben des Datenschutzbeauftragten. \hfill (3~P.)
  \item \textbf{Erläutern} Sie, warum der IT-Leiter in der Regel nicht Datenschutzbeauftragter werden kann. \hfill (3~P.)
  \item \textbf{Ordnen} Sie folgenden Wünschen das passende Betroffenenrecht (mit Artikel) \textbf{zu}: \hfill (5~P.)
    \begin{enumerate}[label=\arabic*., leftmargin=*]
      \item Ein Kunde möchte wissen, welche Daten über ihn gespeichert sind.
      \item Ein Kunde verlangt die Löschung seiner Daten.
      \item Eine falsch geschriebene Adresse soll korrigiert werden.
      \item Ein Kunde will seine Daten in maschinenlesbarer Form erhalten.
    \end{enumerate}
\end{enumerate}

\subsection*{Aufgabe~4 --- Datensicherung (14 Punkte)}

\noindent\textbf{\color{bfwAccent}YouTube-Vorbereitung:}\enspace \href{https://www.youtube.com/watch?v=T2joUFzI_9g}{Backup-Generationenprinzip} (\url{https://www.youtube.com/watch?v=T2joUFzI_9g})

\medskip
\begin{enumerate}[label=(\alph*)]
  \item \textbf{Erläutern} Sie das Generationenprinzip (Großvater--Vater--Sohn) und die Mindestzahl der Generationen. \hfill (4~P.)
  \item \textbf{Ordnen} Sie den Gefahrenbereichen \emph{Datenerfassung, Datenübertragung, Rechnerausfall, Computerkriminalität} je eine Maßnahme \textbf{zu}. \hfill (4~P.)
  \item \textbf{Erklären} Sie die Funktion einer USV. \hfill (3~P.)
  \item \textbf{Unterscheiden} Sie den Aktenvernichter (Datenschutz) vom Backup (Datensicherung). \hfill (3~P.)
\end{enumerate}

\subsection*{Aufgabe~5 --- Umweltgesetze \& Abfallhierarchie (15 Punkte)}

\noindent\textbf{\color{bfwAccent}YouTube-Vorbereitung:}\enspace \href{https://www.youtube.com/watch?v=pV-wTLqZHEQ}{Abfallhierarchie \& KrWG} (\url{https://www.youtube.com/watch?v=pV-wTLqZHEQ})

\medskip
\begin{enumerate}[label=(\alph*)]
  \item \textbf{Nennen} Sie die fünf Stufen der Abfallhierarchie nach \S~6 KrWG in der richtigen Reihenfolge. \hfill (5~P.)
  \item \textbf{Erläutern} Sie den Unterschied zwischen Wiederverwendung und Weiterverwendung mit je einem Beispiel. \hfill (4~P.)
  \item \textbf{Beschreiben} Sie, was das Verpackungsgesetz regelt, und \textbf{nennen} Sie die drei Säulen des dualen Systems. \hfill (4~P.)
  \item \textbf{Erklären} Sie, was das Symbol der durchgestrichenen Mülltonne (ElektroG) bedeutet. \hfill (2~P.)
\end{enumerate}

\subsection*{Aufgabe~6 --- Nachhaltiges Wirtschaften \& Green IT (12 Punkte)}

\noindent\textbf{\color{bfwAccent}YouTube-Vorbereitung:}\enspace \href{https://www.youtube.com/watch?v=zB-njPQZJcQ}{Nachhaltigkeit \& Umweltsiegel} (\url{https://www.youtube.com/watch?v=zB-njPQZJcQ})

\medskip
\begin{enumerate}[label=(\alph*)]
  \item \textbf{Erklären} Sie den Begriff \emph{Nachhaltigkeit} und \textbf{nennen} Sie drei Bestandteile nachhaltigen Wirtschaftens. \hfill (4~P.)
  \item \textbf{Ordnen} Sie den Siegeln \emph{Blauer Engel, Energy Star, FSC} ihre Bedeutung \textbf{zu}. \hfill (3~P.)
  \item \textbf{Nennen} Sie vier Maßnahmen, mit denen ein Büro Energie sparen kann. \hfill (3~P.)
  \item \textbf{Beurteilen} Sie, wie ein DMS zugleich dem Umweltschutz dient (Stichwort Green IT). \hfill (2~P.)
\end{enumerate}

\vspace{1em}
\noindent\textbf{Punkte gesamt: 79}

\newpage

\section{Lösungen}\label{loesungen}

\subsection*{Lösung Aufgabe~1}
\begin{loesung}
\textbf{(a)} Ein Dokumentenmanagementsystem (DMS) wandelt ursprünglich papiergebundene
Schriftstücke in elektronische Dokumente um und verwaltet diese datenbankgestützt; es
archiviert die Dokumente digital und macht sie durchsuchbar.

\textbf{(b)} Vier Vorteile: schneller Zugriff auf Dokumente und Informationen, Wegfall
von Verteilerkopien, reduzierter Dokumentenverlust, geringer Platzbedarf für die
Archivierung, Wegfall monotoner Ablagetätigkeiten, unabhängiger Zugriff auch bei
Bearbeitung.

\textbf{(c)} Revisionssicher gespeicherte Dokumente sind \emph{unveränderbar},
\emph{vollständig}, \emph{nachvollziehbar} (protokolliert) und über die gesamte
Aufbewahrungsfrist \emph{wiederauffindbar} (HGB \S~257, GoBD).

\textbf{(d)} Ein DMS ist die technische Grundlage des papierlosen Büros: Dokumente werden
digital empfangen, bearbeitet und archiviert; Papier, Druck und physische Ablage
entfallen weitgehend, was Platz, Kosten und Ressourcen spart.
\end{loesung}

\subsection*{Lösung Aufgabe~2}
\begin{loesung}
\textbf{(a)} \emph{Datenschutz} schützt natürliche Personen vor der Verletzung ihrer
Persönlichkeitsrechte (Missbrauch personenbezogener Daten). \emph{Datensicherheit}
schützt Hard-, Software und Daten vor Verlust, Zerstörung und Missbrauch durch Unbefugte.

\textbf{(b)} Personenbezogene Daten sind alle Informationen, die sich auf eine
identifizierte oder identifizierbare natürliche Person beziehen (Art.~4 Nr.~1
EU-DSGVO). Beispiele: Name, E-Mail-Adresse, Bankdaten, Kreditkartennummer, IP-Adresse.

\textbf{(c)} EU-Datenschutz-Grundverordnung (EU-DSGVO) und Bundesdatenschutzgesetz
(BDSG).

\textbf{(d)} Die Aussage ist im engeren Sinn ungenau: Eine Firewall ist primär eine
Maßnahme der \emph{Datensicherheit} (Schutz des Netzwerks/der Technik vor unerlaubten
Zugriffen). Sie trägt mittelbar zum Datenschutz bei, indem sie personenbezogene Daten
vor unbefugtem Zugriff bewahrt.
\end{loesung}

\subsection*{Lösung Aufgabe~3}
\begin{loesung}
\textbf{(a)} Ja. Nach \S~38 BDSG ist ein Datenschutzbeauftragter zu ernennen, wenn
mindestens zwanzig Personen ständig mit der automatisierten Verarbeitung personenbezogener
Daten beschäftigt sind. Mit 25 Personen ist die Schwelle überschritten.

\textbf{(b)} Drei Aufgaben (Beispiele): Unterrichtung/Beratung der Unternehmensleitung,
Überwachung der Einhaltung der Datenschutzvorschriften, Sensibilisierung/Schulung der
Mitarbeiter, Zusammenarbeit mit der Aufsichtsbehörde, Beratung bei der
Datenschutz-Folgenabschätzung.

\textbf{(c)} Der IT-Leiter müsste sich sonst in bestimmten Bereichen selbst kontrollieren
(Interessenkonflikt). Der Datenschutzbeauftragte muss unabhängig kontrollieren können.

\textbf{(d)} Zuordnung:
\begin{enumerate}[label=\arabic*., itemsep=0pt]
  \item Auskunftsrecht (Art.~15 EU-DSGVO).
  \item Recht auf Löschung / „Recht auf Vergessenwerden” (Art.~17 EU-DSGVO).
  \item Recht auf Berichtigung (Art.~16 EU-DSGVO).
  \item Recht auf Datenübertragbarkeit (Art.~20 EU-DSGVO).
\end{enumerate}
\end{loesung}

\subsection*{Lösung Aufgabe~4}
\begin{loesung}
\textbf{(a)} Von jedem Datenbestand sind mindestens \emph{drei} aufeinander aufbauende
Generationen vorzuhalten: die älteste Sicherung heißt \emph{Großvater}, die zweitälteste
\emph{Vater}, die jüngste \emph{Sohn}. Mit jeder neuen Sicherung rücken die vorhandenen
Bänder eine Generation weiter.

\textbf{(b)} Zuordnung:
\begin{itemize}[itemsep=0pt]
  \item \emph{Datenerfassung}: Prüfziffernverfahren / Plausibilitätskontrolle.
  \item \emph{Datenübertragung}: Prüfbit, Übertragungsprotokolle, Verschlüsselung.
  \item \emph{Rechnerausfall}: Notstromanlage (USV) / Plattenspiegelung (Duplexing).
  \item \emph{Computerkriminalität}: Zugriffskontrolle/Passwort/PIN, Virenschutzprogramme.
\end{itemize}

\textbf{(c)} Eine USV (unterbrechungsfreie Stromversorgung) wird zwischen Computer und
Steckdose geschaltet und versorgt den Rechner bei Stromausfall so lange weiter, bis die
Daten aus dem Hauptspeicher gesichert sind.

\textbf{(d)} Der \emph{Aktenvernichter} dient dem \emph{Datenschutz}: nicht mehr
benötigte Unterlagen werden datenschutzgerecht vernichtet (Schutz der Personen). Das
\emph{Backup} dient der \emph{Datensicherung}: Daten werden vor Verlust und Zerstörung
geschützt, indem Kopien angelegt werden.
\end{loesung}

\subsection*{Lösung Aufgabe~5}
\begin{loesung}
\textbf{(a)} 1.~Vermeidung, 2.~Vorbereitung zur Wiederverwendung, 3.~Recycling,
4.~sonstige (insbesondere energetische) Verwertung, 5.~Beseitigung. Grundsatz:
Vermeidung vor Verwertung, Verwertung vor Beseitigung.

\textbf{(b)} \emph{Wiederverwendung}: ein bereits gebrauchtes Produkt wird für denselben
Zweck erneut genutzt (z.\,B.\ Glaspfandflasche für Mineralwasser). \emph{Weiterverwendung}:
ein gebrauchtes Produkt wird für einen anderen Zweck genutzt (z.\,B.\ alte Autoreifen als
Schiffsfender).

\textbf{(c)} Das Verpackungsgesetz regelt das Inverkehrbringen, die Rücknahme und die
Verwertung von Verpackungen; Hersteller müssen sich registrieren und an einem dualen
System beteiligen. Drei Säulen: Altpapier, Glas, gelber Sack/gelbe Tonne.

\textbf{(d)} Die durchgestrichene Mülltonne kennzeichnet Elektrogeräte: Sie dürfen
\emph{nicht} über den Restmüll entsorgt, sondern müssen als Elektroschrott gesondert
gesammelt und entsorgt werden.
\end{loesung}

\subsection*{Lösung Aufgabe~6}
\begin{loesung}
\textbf{(a)} \emph{Nachhaltigkeit} (Begriff aus der Forstwirtschaft) bedeutet, heutige
Bedürfnisse so zu befriedigen, dass den nachfolgenden Generationen alle Ressourcen in
gleicher Menge und Qualität zur Verfügung stehen. Bestandteile (drei): umweltfreundliche
Produkte, Reduzierung des Energieverbrauchs, umweltschonende Produktionsverfahren,
rationelle Ressourcenverwendung, fairer Handel.

\textbf{(b)} \emph{Blauer Engel}: Umweltzeichen der Bundesregierung für umwelt- und
gesundheitsfreundliche Produkte. \emph{Energy Star}: Kennzeichnung stromsparender
IT-Geräte (EPA). \emph{FSC}: Holz-/Papierprodukte aus verantwortungsvoll
bewirtschafteten Wäldern.

\textbf{(c)} Vier Maßnahmen (Beispiele): Licht/Bildschirm beim Verlassen abschalten,
Drucker erst bei Bedarf einschalten, stoßlüften statt Kippstellung, bewegungsgesteuertes
Licht, Geräte zum Feierabend komplett abschalten (Stand-by vermeiden), auf LED umstellen,
energieeffiziente Geräte beschaffen.

\textbf{(d)} Ein DMS spart Papier, Druck und physische Archivflächen und unterstützt das
papierlose Büro. Als Teil der \emph{Green IT} schont es Ressourcen über den gesamten
Lebenszyklus --- insbesondere, wenn energieeffiziente Server oder Cloud-Dienste genutzt
werden.
\end{loesung}

\end{document}
